% =============================================================================
\begin{abstract}
\noindent
Este trabalho apresenta a implementação, a partir do zero, de um arcabouço de
Programação Genética (PG) com representação em árvore aplicado à regressão
simbólica em três bases de dados (\code{synth1}, \code{synth2} e
\code{concrete1}). A investigação foi conduzida em duas fases. Na
\textbf{Fase~1}, seguindo o roteiro canônico da disciplina, foram estudados o
tamanho da população e o número de gerações, os métodos de seleção (roleta,
torneio e $\eps$-lexicase), as probabilidades de cruzamento e mutação, o
elitismo e o \emph{bloating}, sempre com 30 repetições independentes e testes
não paramétricos. A $\eps$-lexicase foi o único fator com efeito
estatisticamente significativo no \code{synth1} (RMSE de teste
$5.52 \to 1.62$, $p = 4.1\times10^{-6}$), mas a replicação nas demais bases
expôs três patologias: explosões assintóticas fora da amostra no
\code{synth2} (RMSE máximo de 415), subajuste estrutural no \code{concrete1}
($R^2 \approx 0.45$) e crescimento sistemático das árvores. Na
\textbf{Fase~2}, esses sintomas foram interpretados como propriedades da
\emph{fitness landscape} (polos, rugosidade epistática, perda de diversidade,
múltiplas bacias, platôs neutros e barreiras de escala), originando seis
hipóteses de melhoria (H1--H6) avaliadas por ablação isolada e por
empilhamento sequencial. O \emph{linear scaling} e o \emph{analytic quotient}
eliminaram as explosões do \code{synth2} (30/30 vitórias sobre o estado
anterior) e elevaram o $R^2$ do \code{concrete1} para $0.78$ com árvores
$42\%$ menores; já a combinação de rejeição de duplicatas, parcimônia
lexicográfica e mutação mista reduziu a mediana do erro no \code{synth1} de
$4.28$ (Fase~1) para $0.093$, recuperando a expressão geradora exata em 7 das
30 execuções. Mostra-se ainda que os efeitos são fortemente dependentes da
base (o que ajuda uma paisagem atrapalha outra) e propõe-se a melhor
configuração por base e uma configuração global de compromisso
(\cfg{Stack~B3}), que supera a Fase~1 nas três bases.

\end{abstract}
